\section{Skills 的加载机制}

\subsection{Codex 中的 System Prompt 架构}

Codex 在用户发出第一条请求之前，已经植入了四条系统消息：

\begin{enumerate}
    \item \textbf{System 级提示}：系统级约束。
    \item \textbf{Developer 消息}：权限与沙箱规则。
    \item \textbf{User 消息（注入）}：包含 \texttt{agents.md} 和所有可用 Skills 的描述——这是 Skills 的\textbf{第一次暴露}。
    \item \textbf{Developer 消息}：定义协作模式（Plan mode / Agent mode）。
\end{enumerate}

\subsection{Skill 的匹配与加载流程}

\begin{enumerate}
    \item 用户发出真实请求。
    \item Coding Agent 基于 query \textbf{匹配最合适的 Skill}。
    \item Agent 调用 \texttt{execute\_commands} 工具，执行 \texttt{cat} 命令加载对应 \texttt{skill.md} 的全文。
    \item Skill 内容进入上下文，后续所有交互均基于该 Skill 的描述执行。
    \item 一个 turn 结束后，用户可继续迭代需求。
\end{enumerate}

\begin{knowledgebox}{agents.md 的两级来源}
\begin{itemize}
    \item \textbf{系统级}：\texttt{\~{}/.codex/agents.md}（全局配置）。
    \item \textbf{项目级}：代码仓库根目录下的 \texttt{agents.md}（项目专属指令）。
\end{itemize}
两者都会被注入到第三条 User 消息中。
\end{knowledgebox}

\subsection{安装与更新}

\begin{itemize}
    \item Claude Code 通过 plugin 机制安装 Superpowers。
    \item Codex 通过自然语言让 Agent 自行安装。
    \item 更新方式：进入 Superpowers 目录执行 \texttt{git pull}，下次开启新会话即自动加载。
    \item 当前版本支持\textbf{自主发现}：输入 \texttt{/skills} 可列出所有可用 Skill。
\end{itemize}

\subsection{本章小结}

Skills 通过 Prompt Injection 的方式，在 Codex 的第三条 User 消息中暴露目录；Agent 根据用户 query 动态匹配并用 \texttt{cat} 命令加载具体的 \texttt{skill.md}，从而将软件工程流程注入模型上下文。


